\section{Experimenting with RLRom}
\label{sec:workflow}

A typical study in this thesis follows a short, scriptable loop that keeps
training, model selection and evaluation mutually consistent. Because the
specification, model, and test parameters all live in the same file, the whole loop
is reproducible from a single artifact and trivially adaptable to a new objective by
editing YAML rather than code.

\subsection{Training across seeds}
To make results robust and to estimate variance, the library supports repeating the
same training several times, each run starting from scratch in its own
\texttt{<name>\_<date>\_\_trainingN} folder. This is enabled by a single setting in
the configuration:
\begin{cfglisting}[caption={Repeating training across seeds.},label={lst:repeat}]
cfg_train:
  repeat_training: 5
  # ...other training options
\end{cfglisting}

\subsection{Selecting a model}
After training (possibly several times), a \emph{model selection} step picks the
best checkpoint and \emph{activates} it through symlinks, so that subsequent
testing loads it deterministically:
\begin{quote}\ttfamily
rlr select-model examples/cartpole/cfg0tr\_stl\_go\_left.yml
\end{quote}
Selection can rank checkpoints by the mean episode reward (the default), or by a
\emph{formula} metric --- e.g.\ the satisfaction ratio of \texttt{phi\_left\_goal}
via \texttt{-m phi\_left\_goal/ratio\_init\_sat}. Options let one choose the pool of
trainings (\texttt{--trainings last|all}), take the first checkpoint after a given
step (\texttt{--step}), or inspect a ranked table without selecting
(\texttt{--show-top}).

\subsection{Testing and plotting}
The active model is then evaluated with \texttt{rlr test} while the STL formulas
are monitored (Section~\ref{sec:testing}). Finally, training progress can be plotted
from the checkpoint result files (\texttt{res\_step\_*.yml}) produced by
\texttt{rlr train}:
\begin{quote}\raggedright\ttfamily
rlr plot examples/cartpole/cfg0tr\_stl\_go\_left.yml --envelope -o training.png
\end{quote}
With \texttt{--envelope} the tool overlays the mean and the min--max band across the
repeated trainings, and with \texttt{-m METRIC} (repeatable) it can plot specific
formula metrics --- such as a specification's satisfaction ratio --- alongside the
reward, one subplot per metric. This is the pipeline used to produce the learning
curves reported in the experimental chapters.

\subsection{Command-line interface}
\label{sec:workflow:cli}
Table~\ref{tab:cli} summarizes the commands. Most operations also have a thin
standalone alias (\texttt{rlrom\_test}, \texttt{rlrom\_train},
\texttt{rlrom\_plot}, \texttt{rlrom\_select\_model}) besides the aggregated
\texttt{rlr} entry point.

\begin{table}[htbp]
  \centering
  \caption{Command-line commands of \textsc{RLROM}.}
  \label{tab:cli}
  \small
  \begin{tabular}{@{}llp{9.6cm}@{}}
    \toprule
    \textbf{Command} & & \textbf{Function} \\
    \midrule
    \texttt{rlr show cfg.yml} & & print a configuration, optionally expanded
      (\texttt{--expand}) or restricted to a sub-config (\texttt{--subcfg}) \\
    \texttt{rlr test cfg.yml} & & play episodes with a trained model (or
      \texttt{random}/keyboard) and monitor the STL formulas \\
    \texttt{rlr train cfg.yml} & & train an agent, optionally with STL specs \\
    \texttt{rlr plot cfg.yml} & & plot training curves / envelopes from checkpoint
      result files \\
    \texttt{rlr select-model cfg.yml} & & pick the best checkpoint and activate it
      through symlinks for subsequent testing \\
    \bottomrule
  \end{tabular}
\end{table}
